\chapter{Programmer Documentation}\label{ch:c}

\section{Modules and Classes}

The software is programmed in TypeScript, HTML, and CSS\@.
It is divided into six logical modules.
Each module groups related classes by responsibility and could theoretically be used independently,
though the modules are not separated at the source code level.

\oneimage{documentation/modules}{
    Diagram of all modules and their classes.
    Arrows indicate that one module depends on another.
}{0.8\textwidth}

\subsection*{Graphics}

This module acts as a layer between the \texttt{DOM API} and the rest of the software.
It implements common graphical operations, allowing callers to add or modify visual elements
through a single method call instead of calling multiple \texttt{HTMLElement}s directly.

\begin{description}
    \item[\texttt{Frame}]
    The class representing a whole Sudoku puzzle.
    It creates all instances of \texttt{Cell}, \texttt{Edge}, and \texttt{Canvas} needed for the puzzle.
    It provides coordinate-based access to the instances of \texttt{Cell} and \texttt{Edge}.
    Some changes need to be batched and only applied when the method \texttt{Frame.Update()} is called,
    avoiding the cost of updating after every call.

    \item[\texttt{Cell}]
    The class representing a single cell of a Sudoku puzzle, managing the text and styles inside.
    Cells are accessed by a 2D coordinate system \texttt{(y, x)}.

    \item[\texttt{Edge}]
    The class representing a single label of an edge variant.
    All instances of this class are created by \texttt{Frame} and allow the addition of a single label.
    Edges are accessed by a 3D coordinate system \texttt{(direction, y, x)} where
    \texttt{direction} is either \texttt{0} for edges on right side of a cell or \texttt{1} for edges on left side.

    \item[\texttt{Canvas}]
    The class representing graphical elements spanning multiple cells, rendered onto an \texttt{HTMLCanvasElement}.
    It adds two types of graphical elements:
    lines going through cell centers, used to mark diagonal Sudoku regions;
    and borders around groups of cells, used to highlight extra Sudoku regions.
    Elements are added by specifying a set of cells using either the 2D coordinate system
    of \texttt{Frame} or the 1D coordinate system of \texttt{State}.
    The coordinates are then converted to pixel positions on the \texttt{HTMLCanvasElement}.
\end{description}

\subsection*{Data}

This module stores the data representing a Sudoku puzzle's rules, task, and current solving state.

\begin{description}
    \item[\texttt{Sudoku}]
    The class representing a whole Sudoku puzzle.
    It contains \texttt{Regions},
    two instances of \texttt{State} representing the puzzle task and the current solving state,
    and optionally \texttt{EdgeVariant}.

    \item[\texttt{State}]
    The class representing a Sudoku state.
    For each cell, it stores the set of candidate numbers that can be placed in that cell.
    It also provides few methods to update the state.
    \texttt{State} is designed to contain only data that changes during solving.
    Currently, this is limited to the set of candidate numbers for each cell,
    but the class is intended to accommodate additional solving data in the future.


    \item[\texttt{Regions}]
    The class representing Sudoku regions.
    It stores all regions categorized in five types: rows, columns, irregular, diagonal, and extra.
    It provides convenience methods for adding regions and retrieving them by type.

    \item[\texttt{EdgeVariant}]
    The class representing an Edge variant.
    It stores all the information about the variant rules, graphics, and the current labels on each edge.

    \item[\texttt{EdgeVariantFactory}]
    A static class used to create instances of \texttt{EdgeVariant}.
\end{description}

\subsection*{Strategy}

This module implements all available strategies.
Each strategy accepts a \texttt{State} and returns a \texttt{StrategyResult},
which describes how to visualize the result in the editor and which candidate numbers are eliminated.

\begin{description}
    \item[\texttt{SudokuStrategy}]
    The class implementing strategies for basic Sudoku.

    \item[\texttt{PerfectMatching}]
    The class implementing Algorithm~\ref{alg:algorithm} that is needed by Perfect matching strategies.

    \item[\texttt{EdgeStrategy}]
    The class implementing strategies for Edge variants.

    \item[\texttt{StrategyUtils}]
    A class containing methods that are needed by multiple classes in this module.
\end{description}

\subsection*{Editor}

This module implements the editor.

\begin{description}
    \item[\texttt{MainEditor}]
    The class responsible for the editor menu.
    It handles menu button interactions and tracks which section of the menu is currently open.

    \item[\texttt{RegionEditor}]
    The class responsible for editing Jigsaw and Extra regions.
    It handles user input and graphics while the \texttt{Edit Regions} menu is active,
    covering both the \texttt{Jigsaw} and \texttt{Rectangular} subpanels.

    \item[\texttt{NumberEditor}]
    The class responsible for editing Sudoku numbers.
    It handles user input.

    \item[\texttt{EdgeEditor}]
    The class responsible for editing Edge variant.
    It handles user input and configures the Edge variant.

    \item[\texttt{StrategyEditor}]
    The class responsible for strategies.
    It manages a list of all available strategies, runs them on demand, and renders the strategy result.

    \item[\texttt{EditorUtils}]
    A class containing methods that are needed by multiple classes in this module.
\end{description}

\subsection*{Renderer}

This module only contains one class.
It acts as a bridge between the modules \texttt{Data} and \texttt{Graphics}.
User can configure a sudoku using the classes from the module \texttt{Data} and then

\begin{description}
    \item[\texttt{Renderer}]
    The class responsible for drawing a whole Sudoku from a \texttt{State} into HTML using a single method call.
\end{description}

\subsection*{Utils}

This module contains all methods that can be used by multiple modules.

\begin{description}
    \item[\texttt{ArrayUtils}]
    Useful methods for working with arrays.

    \item[\texttt{BinaryUtils}]
    Useful methods for working with binary representation of candidates.
    See Section~\ref{sec:data}.

    \item[\texttt{GraphicsUtils}]
    Useful methods for working with HTML\@.

    \item[\texttt{PositionUtils}]
    Useful methods for working with 1D and 2D coordinate systems.
    See Section~\ref{sec:data}.
\end{description}

\subsection*{Combination of Modules}

The software uses all modules, but it is possible to use some combinations of the modules without the rest.
Module \texttt{Graphics} alone is sufficient to render grid-based puzzles into HTML\@.
Modules \texttt{Graphics}, \texttt{Renderer}, and \texttt{Data} together allow
configuring and rendering a full Sudoku task or solution.
Modules \texttt{Data} and \texttt{Strategy} together allow configuring a Sudoku puzzle and applying solving strategies.

\section{Representation of Data}\label{sec:data}

\subsection*{Binary Representation of Candidates}

The class \texttt{State} stores all candidate numbers in a Sudoku state.
Each cell is encoded as a 31-bit bitset stored as a JavaScript \texttt{number},
where a Sudoku number $m$ is a candidate if and only if \texttt{1 << ($m$ - 1)} is a 1-bit.
The most significant bit is always 0 to prevent the number from being negative.
This binary representation limits the Sudoku size to $31$, though the editor enforces a stricter limit of $16$.

\subsection*{Coordinate Systems}

\subsubsection*{Cell}

The module \texttt{Graphics} uses a 2D coordinate system \texttt{(y, x)} for cells in the HTML \texttt{<table>},
where \texttt{y = 0} is the top row and \texttt{x = 0} is the left-most column.
The module \texttt{Data} uses a 1D coordinate system \texttt{index = y * Sudoku.size + x} for Sudoku cells.
The class \texttt{Renderer} decides the mapping of \texttt{State} into \texttt{Frame}.

\subsubsection*{Edge}

Edge variant labels use a 3D coordinate system \texttt{(direction, y, x)},
where \texttt{direction} is either \texttt{0} for the right side of a cell or \texttt{1} for the bottom.

\section{Extensibility}\label{sec:extensibility}

\subsection*{Variants}

To add a new variant, add a field to \texttt{Sudoku} that represents the variant's data
and implement a solving strategy for it.
To add the necessary UI to the editor, add the needed buttons to \texttt{editor.html}
and implement their event handlers in \texttt{MainEditor}.
To add new graphical elements, implement needed CSS styles to file \texttt{code/src/styles/puzzle.css},
it might be necessary to update \texttt{Renderer} or the module \texttt{Graphics}
if the graphical elements are too different.

\subsection*{Edge variants}
A new Edge variant strategy can be added by appending an element to \texttt{EdgeVariantFactory.variants}.
Each element should be created using the \texttt{EdgeVariantFactory} constructor,
which requires four parameters.
The last step is to implement all CSS classes specified in the third parameter.
The following code shows an example of adding the Edge variant VX\@.

\noindent
\begin{minipage}{\textwidth}
\begin{lstlisting}[language=JavaScript]
// call EdgeVariantFactory constructor and fill the parameters
vx: new EdgeVariantFactory(
    // name of the variant shown on the button in editor
    "VX",
    // function that returns all possible labels
    // for combination of two numbers in the edge-adjacent cells
    // size is Sudoku.size - the size of the Sudoku puzzle
    // value1 is the number inside the top or the left cell
    // value2 is the number inside the bottom or the right cell
    // output is an array with all labels
    // that can appear between the two cells
    (size: number, value1: number, value2: number): number[] => {
        switch (value1 + value2) {
            // if the sum is 5, return label V
            case 5:     return [1];
            // if the sum is 10, return label X
            case 10:    return [2];
            // otherwise return label with no graphical element
            default:    return [0];
        }
    },
    // return array of label graphics
    (size: number): (EdgeLabelStyle | null)[] => [
        null, // no graphical element
        { css: ["vx-v"] }, // CSS class for V
        { css: ["vx-x"] }, // CSS class for X
    ],
    // is the variant applicable for this size of sudoku?
    (size: number): boolean => true,
),
\end{lstlisting}
\end{minipage}

\texttt{EdgeEditor} iterates over all elements of \texttt{EdgeVariantFactory.variants}
and generates the corresponding buttons for adding Edge variants in the editor.
When an Edge variant is added, \texttt{EdgeVariantFactory.build} converts the data into
the format expected by the \texttt{EdgeVariant} constructor,
and \texttt{EdgeEditor} adds it to \texttt{Sudoku}.

\subsection*{Solving Strategies}

A strategy can be added by appending an element of type \texttt{StrategyTypeData}
to the array \texttt{StrategyEditor.strategyTypeData}.
The strategies are in this array in the same order in which the \texttt{StrategyEditor} tries to apply them,
when the user presses the \editorbutton{Apply Step}.
The first strategy that returns a \texttt{StrategyResult} with \texttt{isUpdated = true} will be applied.
The following code shows an example of adding the strategy Edge Region.

\noindent
\begin{minipage}{\textwidth}
\begin{lstlisting}[language=JavaScript]
{
    // name of the strategy on the button in the editor
    name: "Edge Region",
    // method implementing the strategy
    // it has one parameter State and returns StrategyResult
    // see comment on StrategyResult for more information
    strategy: EdgeStrategy.region.bind(EdgeStrategy),
    // is the strategy applicable on the sudoku?
    isUsable: (sudoku: Sudoku) => sudoku.hasEdgeVariant(),
    // can strategy proof that Sudoku has no solution?
    isValidator: false,
},
\end{lstlisting}
\end{minipage}

\subsection*{Bigger Sudoku}

The size of a Sudoku puzzle is limited to $16$ by \texttt{MainEditor.MAX\_SIZE}.
The main reason is that the number of candidates in a Sudoku state is $n^3$,
as there are $n^2$ cells each with up to $n$ candidate numbers.
To add support for larger Sudoku puzzles,
it is better to avoid candidate numbers and use the strategies described in Section~\ref{sec:paper}.
Additionally, larger Sudoku puzzles are generally too complex for humans to solve
using strategies that require candidate numbers,
making candidate-free strategies more practical for puzzles beyond size $16 \times 16$.
